\documentclass[11pt]{article}
\usepackage[margin=1.1in]{geometry}
\usepackage{amsmath,amssymb,amsthm}
\usepackage{xcolor}
\usepackage[colorlinks=true,linkcolor=blue!50!black,citecolor=blue!50!black,urlcolor=blue!50!black]{hyperref}

\newtheorem{theorem}{Theorem}[section]
\newtheorem{lemma}[theorem]{Lemma}
\theoremstyle{remark}
\newtheorem{remark}[theorem]{Remark}


\newcommand{\Z}{\mathbb{Z}}
\newcommand{\R}{\mathbb{R}}
\newcommand{\C}{\mathbb{C}}
\newcommand{\lk}{\operatorname{lk}}

\title{Zippers: an informal note\\on the mathematics behind the programs}
\author{Danny Calegari \and Claude (Anthropic)}
\date{October 2026}

\begin{document}
\maketitle

\begin{center}
\fbox{\parbox{0.88\textwidth}{\small This is an informal note, \emph{not for publication}. It describes some elements of
the mathematical theory of zippers, as the programs of the \texttt{zippers} suite illustrate them. Theorems~\ref{thm:conn} and~\ref{thm:link} are
work in progress: what is given for them is a sketch of a proof, not a proof.}}
\end{center}

\section{Zippers}

Let $M$ be a closed manifold with word-hyperbolic fundamental group $G$, and let $\varphi\colon G\to\Z$ be a
surjective homomorphism, for instance the one induced by a map $M\to S^1$. Write $K=\ker\varphi$. The universal cover
$\widetilde M$ is $\delta$-hyperbolic, and $G$ acts on it and on its Gromov boundary $\partial_\infty G$. The map $M\to S^1$ lifts to a function $f\colon\widetilde M\to\R$ with $f(gx)=f(x)+\varphi(g)$: a height.

A \emph{staircase} is a ray in $\widetilde M$ that goes to infinity at a linear speed, and along which the height $f$
grows at a linear rate. (Precisely: a quasigeodesic ray along which $f$ increases linearly.) Every staircase converges to a
point of $\partial_\infty G$. Following Calegari and Loukidou~\cite{CL}:
\[
Z^+=\{\text{the endpoints of staircases}\},\qquad Z^-=\{\text{the endpoints of staircases for } -f\}.
\]

\begin{theorem}[Calegari--Loukidou \cite{CL}]\label{thm:CL}
If $K$ is finitely generated, then $Z^+$ and $Z^-$ are nonempty, disjoint, path-connected and invariant under
$G$; in particular both are dense in $\partial_\infty G$.
\end{theorem}

They are disjoint for a simple reason: two rays with the same endpoint stay a bounded distance apart, so they cannot have
heights going to $+\infty$ on one and $-\infty$ on the other.

\begin{remark}\label{rem:quasi}
Theorem~\ref{thm:CL} holds more generally whenever $\varphi\colon G\to\Z$ is a quasimorphism whose coarse level sets are coarsely
connected. In this note we consider only homomorphisms, for simplicity.
\end{remark}

\begin{remark}\label{rem:cusps}
If $M$ has cusps, so that $G$ is relatively hyperbolic, the sets $Z^+$ and $Z^-$ can still be constructed, but they are no
longer disjoint.
\end{remark}

\section{How the programs draw them: base arcs from a presentation}

\paragraph{The input.} The programs start from a \emph{presentation} of the situation:
\begin{itemize}
\item a finite set $S\subset K$ that generates $K$ and contains the inverse of each of its elements (the \emph{letters});
\item an element $t\in G$ with $\varphi(t)>0$;
\item for each letter $s$, its \emph{roof}: an expression
\[
t^{-1}st=x_1x_2\cdots x_k\qquad(x_i\in S).
\]
It exists because $K$ is normal: $t^{-1}st$ lies in $K$, so it is a word in the letters.
\end{itemize}

\paragraph{The base arcs.} Let $V$ be the attracting fixed point of $t$ on $\partial_\infty G$. For each letter $s$ there is an arc
$\alpha_s$ in $\partial_\infty G$ from $V$ to $sV$, determined by the self-similarity
\[
\alpha_s \;=\; t\cdot\bigl(\alpha_{x_1}\,\cup\, x_1\alpha_{x_2}\,\cup\, x_1x_2\alpha_{x_3}\,\cup\cdots\cup\,
x_1\cdots x_{k-1}\alpha_{x_k}\bigr).
\]
The right-hand side is a path: $x_1\cdots x_{i-1}\alpha_{x_i}$ runs from $x_1\cdots x_{i-1}V$ to $x_1\cdots x_iV$, and
applying $t$ (which fixes $V$) gives a path from $V$ to $tx_1\cdots x_kV=stV=sV$. The arcs $\alpha_s$ are the
\emph{base arcs}. Every point of a base arc is the endpoint of a staircase, so the union of the translates
$g\alpha_s$, $g\in G$, is a path-connected, $G$-invariant subset of $Z^+$. The same construction with $t^{-1}$,
its attracting fixed point, and the \emph{down roofs} $tst^{-1}=y_1\cdots y_l$ gives the base arcs of $Z^-$.
\begin{remark}\label{rem:union}
The union of the translates of the base arcs is contained in $Z^+$. It is probably equal to $Z^+$, but we have not
thought about this carefully.
\end{remark}

\paragraph{What the programs compute.} A translate $h\alpha_s$ (a \emph{branch}) is replaced by the $k$ branches on the
right-hand side, $ht\alpha_{x_1}, htx_1\alpha_{x_2},\dots$, until the chord joining the two ends of each branch is shorter
than the accuracy (the option \texttt{-a}); such a branch is drawn as a segment. All group elements are kept exactly, and
a branch reached twice is drawn once. The picture is the union of the translates $g\alpha_s$ for the elements $g$ of a
ball of $G$ (those with $d(o,go)\le R$, the option \texttt{-R}), nearest first. With \texttt{-R 0} it is the base arcs
alone.

\section{Roofs}

Arrange the elements of $G$ by their height $\varphi$. The letter $s$ is a step from $1$ to $s$ at height $0$. Its roof
is a path one story up, from $t$ to $st$, through the points
\[
t,\quad tx_1,\quad tx_1x_2,\quad\dots,\quad tx_1\cdots x_k=st,
\]
all at height $\varphi(t)$. Each step of the roof is again a letter, so it has a roof of its own, one story higher, and so on.
These nested roofs are the \emph{blocks} of~\cite{CL}; a staircase climbs through them, and the base arc $\alpha_s$ is the set of
the endpoints of the staircases that start on the step from $1$ to $s$.

The presentation is not unique: $S$, $t$ and the roofs can all be chosen in many ways. Different choices give different
base arcs, but $Z^+$ and $Z^-$ depend only on $G$ and $\varphi$. The choice matters for the pictures:
\begin{itemize}
\item the kernel search of the suite (\texttt{zippers search}) finds $S$, $t$ and the roofs, and checks every roof exactly;
\item it can choose short roofs, and roofs that make the two orientations of an arc coincide; with them, far less of the
picture is drawn twice (the overdraw of $Z^-$ in $S^4$ falls from $79\%$ to $5\%$; \texttt{docs/developer\_guide.md}).
\end{itemize}

\section{Coarse connectivity of the level sets, and connectivity at infinity}

A \emph{level set} is $L_c=f^{-1}(c)\subset\widetilde M$. It is \emph{coarsely $n$-connected} if there is a constant $D$
such that every map of a sphere $S^i$, $i\le n$, into $L_c$ extends to a map of the ball $B^{i+1}$ into the
$D$-neighbourhood of $L_c$. For $n=0$ this says that $L_c$ is coarsely connected, which holds exactly when $K$ is finitely
generated; in general, the level sets are coarsely $n$-connected exactly when $K$ is of type $F_{n+1}$ (Gromov~\cite{G},
where coarse higher connectivity is introduced).

\begin{theorem}[Calegari, work in progress]\label{thm:conn}
If the level sets of $f$ are coarsely $n$-connected, then $Z^+$ and $Z^-$ are $n$-connected: every map of a sphere $S^i$,
$i\le n$, into $Z^\pm$ extends to a map of the ball $B^{i+1}$ into $Z^\pm$.
\end{theorem}

The key to the first step of the proof is a lemma about the geodesics that run from a point of $Z^-$ to the points of
a compact subset of $Z^+$, in its weak topology. Fix a basepoint $o$. For $K\ge1$ let $Z^+(K)$ be the set of the endpoints
of the geodesic rays that start within distance $K$ of $o$ and along which the height tends to $+\infty$ and is a
$K$-quasi-isometry to $\R$; define $Z^-(K)$ in the same way with $-f$. Then $Z^+$ is the increasing union of the sets
$Z^+(K)$, and it has the \emph{weak topology}: a set is open when its intersection with each $Z^+(K)$ is open (in the
topology of $\partial_\infty G$). The sets $Z^+(K)$ are not invariant under $G$, but $Z^+$ is, and $G$ acts on it by
homeomorphisms of the weak topology. A compact subset of $Z^+$, such as the image of a sphere or of a null-homotopy of
one, lies in a single $Z^+(K)$.

\begin{lemma}\label{lem:geod}
Let $X$ be a compact subset of $Z^+$, and let $p$ be a point of $Z^-$. Let $\Gamma(X)$ be the family of the geodesics from $p$
to the points of $X$. Then the height $f$ is a quasi-isometry to $\R$ along every geodesic $\gamma$ in $\Gamma(X)$, with
constants that do not depend on $\gamma$.
\end{lemma}

\begin{proof}
Since $X$ is compact, it lies in a single $Z^+(K)$, and $p$ lies in some $Z^-(K')$. For $x\in X$ let $r_x$ be a ray from
near $o$ to $x$ along which the height is a $K$-quasi-isometry, and $r^-$ a ray from near $o$ to $p$ along which $-f$ is a
$K'$-quasi-isometry.
\begin{enumerate}
\item The union $r^-\cup r_x$ is \emph{expanding}: any two of its points $y,y'$ satisfy $|f(y)-f(y')|\ge d(y,y')/A-B$,
with $A$, $B$ depending only on $K$, $K'$ and the Lipschitz constant of $f$. On each ray this is the quasi-isometry. For
$y\in r^-$ and $y'\in r_x$ it follows from $f(y)\le f(o)-d(o,y)/A+B$, $f(y')\ge f(o)+d(o,y')/A-B$ and the triangle
inequality.
\item The geodesic $\gamma$ from $p$ to $x$ stays within a bounded distance of $r^-\cup r_x$, since ideal polygons in a
hyperbolic space are thin and the two rays start near $o$.
\item Being expanding passes to bounded neighbourhoods, since $f$ is Lipschitz. So $|f(\gamma(u))-f(\gamma(u'))|\ge
|u-u'|/A'-B'$, and $f\circ\gamma$ runs from $-\infty$ to $+\infty$. A continuous function on $\R$ with these two properties
increases coarsely: it is a quasi-isometry, with constants that do not depend on $x$.\qedhere
\end{enumerate}
\end{proof}

\noindent\emph{Sketch of a proof of Theorem~\ref{thm:conn}.} This is a sketch only, and the details have not been written out.
\begin{enumerate}
\item Let $S^i\to Z^+$ be a map, with image $X$, and fix a point $p$ of $Z^-$. By Lemma~\ref{lem:geod}, along each geodesic from $p$
to a point of $X$ the height goes from $-\infty$ to $+\infty$ at a uniformly linear rate, so the upper half of each is a
staircase, with uniform constants. These geodesics vary continuously with their endpoints, and each crosses a given
level set in a set of bounded size: the crossings form a coarse sphere in that level set.
\item Since the level set is coarsely $n$-connected, that coarse sphere bounds a coarse ball at the same level.
\item Push the ball up, story by story: each small cell at one level is spanned by a cell one story higher, as each step is
spanned by its roof. This is the higher-dimensional form of the block filling that proves Theorem~\ref{thm:CL} (where $i=0$, and a
path of steps is filled by its roofs).
\item The pushed-up cells form a ball's worth of staircases, and their endpoints give a ball in $Z^+$ bounding the
sphere.
\end{enumerate}
The work lies in step~3: keeping the constants uniform, so that the pushed-up family is still a family of staircases, and
the endpoint map is continuous.

\section{Examples}

\subsection{Hyperbolic 3-manifolds that fiber over the circle: zippers on $S^2$}

The figure-eight knot complement fibers over the circle with fiber a once-punctured torus, whose fundamental group is the
free group $F=\langle a,b\rangle$. The monodromy is the automorphism
\[
\varphi(a)=ab,\qquad \varphi(b)=bab,
\]
the product of the two Dehn twists $R\colon(a,b)\mapsto(a,ba)$ and $L\colon(a,b)\mapsto(ab,b)$, written $RL$ (applied
from left to right). The group is $G=\langle a,b,t\mid t^{-1}gt=\varphi(g)\text{ for } g\in F\rangle$, the
height of an element is its exponent of $t$, and $K=F$. So the roofs are just the monodromy:
\[
t^{-1}at=ab,\qquad t^{-1}bt=bab,
\]
and the two base arcs of $Z^+$ satisfy
\[
\alpha_a=t\,(\alpha_a\cup a\,\alpha_b),\qquad \alpha_b=t\,(\alpha_b\cup b\,\alpha_a\cup ba\,\alpha_b).
\]
\begin{itemize}
\item \texttt{zippers draw fig8 -n inf} draws them for the knot complement, which is cusped (of finite volume, but not
closed).
\item \texttt{zippers draw fig8 -n N} draws them for the closed orbifold obtained by replacing the puncture with a cone point
of order $N$. It is hyperbolic for every $N\ge2$: the fiber is then a hyperbolic $2$-orbifold, a torus with one cone
point, and the monodromy is pseudo-Anosov. The program finds its holonomy by following the fixed point of the trace map
from the cusp.
\item Any positive word in $R$ and $L$ with both letters gives another example (\texttt{zippers draw RRL}).
\item \texttt{-C} also draws the \emph{CaTherine wheel}~\cite{CLwheels}.
\end{itemize}

\subsection{The examples of Llosa Isenrich and Py: zippers on $S^3$}

Llosa Isenrich and Py~\cite{LIP} construct compact complex hyperbolic manifolds $X=G\backslash B^m$ of complex
dimension $m$, with a map to the circle whose height function $f$ is Morse, all of whose critical points have index $m$
(the middle dimension). For them $K$ is of type $F_{m-1}$ but not of type $FP_m$. Here $B^m$ is complex hyperbolic space,
and $\partial_\infty G=S^{2m-1}$.

The simplest is the \emph{Cartwright--Steger surface}~\cite{CS}, with $m=2$: a compact complex hyperbolic surface whose
first Betti number is $2$. Its map to the circle comes from its Albanese map to an elliptic curve, a Lefschetz
fibration with exactly three nodes whose fibers have genus $19$ (see Szab\'o~\cite{Sz} for the Albanese map and the
fundamental group), and $f$ has $3$ critical points in $X$, all of index $2$. So $K$ is finitely generated but not finitely presented, and the zippers are
path-connected subsets of $S^3$ (Theorem~\ref{thm:CL}).
\begin{itemize}
\item \texttt{zippers draw cs} draws them. The default presentation has $36$ letters and an element $t$ that is a word of
length $17$ in the three generators of the lattice.
\item For instance, letter $1$ is $a_1^{-1}a_3^{-1}$, and its roof is $t^{-1}s_1t=s_{32}\,s_8\,s_{19}$. The $36$ roofs have
$6$ letters on average, and at most $14$.
\end{itemize}

\subsection{A near example: a cusped hyperbolic 5-manifold, zippers on $S^4$}

Italiano, Martelli and Migliorini~\cite{IMM} found a hyperbolic $5$-manifold that fibers over the circle. It has finite
volume but is not closed: it has cusps, so $\pi_1$ is relatively hyperbolic and $Z^+$ and $Z^-$ are not disjoint; see Remark~\ref{rem:cusps}.
The programs draw its arcs by the same construction (\texttt{zippers draw imm}). Near a cusp, an arc runs into the cusp's
point at infinity, and the programs draw it through those points. The pictures show a close family resemblance to the
zippers of the figure-eight knot complement, especially near the cusps.

\section{Linking at infinity and homology in the middle dimension}

Let $X=G\backslash B^m$ be one of the examples of Llosa Isenrich and Py. Every critical point $p$ of the height $f$ on
$B^m$ (there are finitely many up to $G$) has:
\begin{itemize}
\item an \emph{ascending thimble} $T^+_p$, the points that flow up from $p$ along the gradient of $f$; it is an
$m$-dimensional plane, since the index of $p$ is $m$;
\item a \emph{descending thimble} $T^-_p$, the points that flow down from $p$; also an $m$-plane.
\end{itemize}

\begin{theorem}[work in progress]\label{thm:link}
The closure of $T^+_p$ in the closed ball $B^m\cup S^{2m-1}$ meets the sphere at infinity in an $(m-1)$-sphere
$v^+_p\subset Z^+$; likewise $T^-_p$ gives an $(m-1)$-sphere $v^-_p\subset Z^-$. Their linking numbers in $S^{2m-1}$ are
\[
\lk(v^+_p,v^-_q)=\delta_{pq}
\]
(with suitable orientations). Consequently $H_{m-1}(Z^\pm)\neq0$, and with Theorem~\ref{thm:conn}, $Z^\pm$ is $(m-2)$-connected but not
$(m-1)$-connected.
\end{theorem}

So the connectivity of the zippers records exactly how finite $K$ is: of type $F_{m-1}$ (Theorem~\ref{thm:conn} gives
$(m-2)$-connected zippers), but not $FP_m$ (the linked spheres).

\noindent\emph{Sketch of a proof.} Again a sketch only.
\begin{enumerate}
\item The gradient lines in $T^+_p$ that start at $p$ are staircases: along them $f$ increases linearly, and they are
quasigeodesic, with constants that do not depend on $p$. So their endpoints lie in $Z^+$, and they make up an
$(m-1)$-sphere at infinity.
\item $T^+_p$ and $T^-_q$ are disjoint when $p\ne q$, and when $p=q$ they meet only at $p$, transversally. So the
intersection number of the two disks in the $2m$-ball is $\delta_{pq}$, and the linking number of their boundaries is the
same.
\item A sphere in $Z^+$ that links a sphere of $Z^-$ cannot be filled in the complement of the second, which contains
$Z^+$. So it is essential in $Z^+$.
\end{enumerate}

\paragraph{What the programs show.} For the Cartwright--Steger surface ($m=2$) the theorem predicts linked loops. The
thimble loops $v^\pm_p$ need not lie in the drawn zippers, but the program \texttt{link3} finds loops in the drawn $Z^+$ and
$Z^-$, made from relations of $K$, and computes their linking numbers: the $13$ pairs of \texttt{data/s3/linked\_pairs.txt}
have linking number $\pm1$. One of them (pair $5$, linking number $-1$) is the $S^3$ picture of the gallery:
\texttt{zippers draw cs --pairs linked\_pairs.txt --pairindex 5 --context}.

\begin{thebibliography}{9}
\bibitem{CL} D.~Calegari and I.~Loukidou, \emph{Zippers}, Geom.\ Topol.\ \textbf{30} (2026), 1931--1962; arXiv:2411.15610.
\bibitem{CLwheels} D.~Calegari and I.~Loukidou, \emph{CaTherine wheels}, arXiv:2604.24619.
\bibitem{CS} D.~I.~Cartwright and T.~Steger, \emph{Enumeration of the 50 fake projective planes}, C.~R.~Math.~Acad.~Sci.\
Paris \textbf{348} (2010), 11--13.
\bibitem{G} M.~Gromov, \emph{Asymptotic invariants of infinite groups}, in \emph{Geometric Group Theory}, Vol.~2
(Sussex, 1991), London Math.\ Soc.\ Lecture Note Ser.\ \textbf{182}, Cambridge Univ.\ Press, 1993, 1--295.
\bibitem{IMM} G.~Italiano, B.~Martelli and M.~Migliorini, \emph{Hyperbolic 5-manifolds that fiber over $S^1$}, Invent.\
Math.\ \textbf{231} (2023), 1--38.
\bibitem{LIP} C.~Llosa Isenrich and P.~Py, \emph{Subgroups of hyperbolic groups, finiteness properties and complex
hyperbolic lattices}, Invent.\ Math.\ \textbf{235} (2024), 233--254; arXiv:2204.05788.
\bibitem{Sz} Z.~Szab\'o, \emph{A polynomial formula for the Albanese map of the Cartwright--Steger surface},
arXiv:2609.23124.
\end{thebibliography}

\end{document}
